%%%%%%%%%%%%%%%%%%%%%%%%%%%%%%%%%%%%%%%%%%%%%%%%%%%%%%%%%%%%%%%%%%%%%%%%%%%%%%%%
% abstract.tex: Abstract
%%%%%%%%%%%%%%%%%%%%%%%%%%%%%%%%%%%%%%%%%%%%%%%%%%%%%%%%%%%%%%%%%%%%%%%%%%%%%%%%

Since 2015, the Laser Interferometer Gravitational-wave Observatory (LIGO) has detected hundreds of signals from merging black holes and neutron stars.
These detections represent the early stages of the field of gravitational-wave astronomy, giving us new tools with which to probe our Universe.
Of particular interest is the advent of multi-messenger astronomy.
In August 2017, after the first observation of gravitational waves from merging neutron stars, and the coincident observation of a kilonova in the electromagnetic spectrum, there has been substantial effort to study these multiply-constrained events.
A repeated merger of this type has not materialized, emphasizing the serendipity of the original observation and the need to be prepared for the next occurrence.

In this dissertation, I present work motivated by the scarcity of multi-messenger events.
I introduce two machine learning methods, \aframe and \amplfi, for the detection and characterization of compact binary coalescences at extremely low-latency.
These models are capable of producing alerts more quickly than their traditional counterparts currently employed by the LVK for its real-time alerts during observation runs.
The underlying software infrastructure, including GPU-accelerated data processing libraries and a modular deployment architecture, is designed to support extension to new source classes and future detector networks.
I discuss the operation of these algorithms and report their performance on real observing data, including archival data from LIGO's third and fourth observing runs as well as live data during the final segment of the fourth observing run.

As the network of ground- and space-based detectors grows with the addition of next-generation instruments, the range of potentially detectable targets expands and the computational cost to cover the full parameter space with traditional methods increases.
Machine learning tools offer a scalable solution to enable future scientific discovery.
The computational savings for low-latency searches, coupled with the innate flexibility of neural networks, make them a natural approach for the upcoming era of gravitational-wave astronomy.
The work presented in this dissertation establishes a foundation upon which novel algorithms can be developed and made ready for production deployment.



%%%%%%%%%%%%%%%%%%%%%%%%%%%%%%%%%%%%%%%%%%%%%%%%%%%%%%%%%%%%%%%%%%%%%%%%%%%%%%%%
